\documentclass[11pt,a4paper]{article}
\usepackage[utf8]{inputenc}
\usepackage[vietnamese]{babel}
\usepackage{amsmath,amssymb}
\usepackage{geometry}
\geometry{a4paper, total={180mm,260mm}, left=15mm, right=15mm, top=18mm, bottom=18mm}
\usepackage{tikz}
\usetikzlibrary{shapes.geometric, arrows.meta, positioning, calc}
\usepackage{tcolorbox}
\tcbuselibrary{skins,breakable}
\usepackage{xcolor}
\usepackage{multicol}
\usepackage{array}
\usepackage{fontawesome}
\usepackage{fancyhdr}

\pagestyle{fancy}
\fancyhf{}
\renewcommand{\headrulewidth}{0.5pt}
\renewcommand{\footrulewidth}{0.5pt}
\fancyhead[L]{\small\bfseries\color{teal}\faBook\ SỔ TAY KIẾN THỨC TOÁN 3}
\fancyhead[R]{\small\bfseries\color{teal}Lớp Toán Thầy Lĩnh -- loptoanthaylinh.vn}
\fancyfoot[C]{\small\color{gray}-- Sổ tay kiến thức: Bài 45. Các số có bốn chữ số. Số 10 000 --}

\definecolor{primary}{HTML}{005F73}
\definecolor{secondary}{HTML}{0A9396}
\definecolor{accent}{HTML}{EE9B00}
\definecolor{danger}{HTML}{AE2012}
\definecolor{boxbg}{HTML}{F4F6F8}

\begin{document}

% --- BANNER TIÊU ĐỀ ---
\begin{tcolorbox}[
  enhanced,
  colback=white,
  colframe=primary,
  arc=3mm,
  boxrule=1.2pt,
  left=4mm, right=4mm, top=2.5mm, bottom=2.5mm,
  drop shadow=black!10!white
]
  \centering
  {\Large\bfseries\color{primary}\faBookmark\ BÀI 45: CÁC SỐ CÓ BỐN CHỮ SỐ. SỐ 10 000}\\[1.5mm]
  {\small\itshape\color{black!75}Cấu tạo thập phân: Nghìn, Trăm, Chục, Đơn vị; đọc viết số và số liền trước, số liền sau trên 1 trang A4}
\end{tcolorbox}

\vspace{-2mm}


\noindent
\begin{minipage}[t]{0.49\textwidth}\vspace{0pt}
\begin{tcolorbox}[
  enhanced,
  colback=blue!4!white,
  colframe=blue!75!black,
  arc=2mm, boxrule=0.8pt,
  left=2mm, right=2mm, top=1.8mm, bottom=1.8mm,
  title={\footnotesize\bfseries\color{white}\faKey\ 1. Cấu tạo số có 4 chữ số}
]
{\footnotesize
\begin{itemize}\setlength{\itemsep}{1.5pt}\setlength{\leftskip}{-2.5mm}
  \item Số có 4 chữ số gồm các hàng tính từ phải sang trái:
  $$\mathbf{\text{Đơn vị}} \to \mathbf{\text{Chục}} \to \mathbf{\text{Trăm}} \to \mathbf{\text{Nghìn}}$$
  \item Phân tích cấu tạo số tổng quát:
  $$\overline{abcd} = a \times 1000 + b \times 100 + c \times 10 + d$$
  \item \textbf{Số 10 000 (mười nghìn / một vạn):}
  Là số nhỏ nhất có $5$ chữ số, liền sau của số $9999$:
  $$9999 + 1 = \mathbf{10\,000}$$
\end{itemize}
}
\end{tcolorbox}
\end{minipage}
\hfill
\begin{minipage}[t]{0.49\textwidth}\vspace{0pt}
\begin{tcolorbox}[
  enhanced,
  colback=teal!4!white,
  colframe=teal!75!black,
  arc=2mm, boxrule=0.8pt,
  left=2mm, right=2mm, top=1.8mm, bottom=1.8mm,
  title={\footnotesize\bfseries\color{white}\faSliders\ 2. Quy tắc đọc và viết số}
]
{\footnotesize
\begin{itemize}\setlength{\itemsep}{1.5pt}\setlength{\leftskip}{-2.5mm}
  \item Đọc và viết số lần lượt từ hàng cao đến hàng thấp (từ trái sang phải).
  \item Chữ số hàng trăm là $0$: đọc là \textbf{"không trăm"}.
  \item Chữ số hàng chục là $0$: đọc là \textbf{"linh"} (hoặc \textbf{"lẻ"}).
  \item Chữ số hàng đơn vị là $5$:
  \begin{itemize}\setlength{\itemsep}{0.5pt}\setlength{\leftskip}{-2.5mm}
    \item Đọc là "lăm" nếu chữ số hàng chục $> 0$.
    \item Đọc là "năm" nếu chữ số hàng chục $= 0$.
  \end{itemize}
\end{itemize}
}
\end{tcolorbox}
\end{minipage}

\vspace{-2.5mm}
% --- KHỐI 3: SƠ ĐỒ TƯ DUY TIKZ ---
\noindent
\begin{tcolorbox}[
  enhanced,
  colback=orange!4!white,
  colframe=orange!80!black,
  arc=2mm, boxrule=0.8pt,
  left=2mm, right=2mm, top=1.5mm, bottom=1.5mm,
  title={\footnotesize\bfseries\color{white}\faSitemap\ SƠ ĐỒ TƯ DUY: HỆ THỐNG CÁC HÀNG SỐ CÓ 4 CHỮ SỐ}
]
\centering
\begin{tikzpicture}[
  u/.style={draw=blue!80!black, fill=blue!10, rounded corners=2pt, font=\scriptsize\bfseries, minimum width=2.4cm, minimum height=0.7cm, align=center},
  arr/.style={-{Stealth[length=2mm]}, thick}
]
  \node[u, fill=purple!20, draw=purple!80!black] (h_nghin) at (0, 0) {HÀNG NGHÌN\\$a \times 1000$};
  \node[u, fill=teal!20, draw=teal!80!black] (h_tram) at (3.5, 0) {HÀNG TRĂM\\$b \times 100$};
  \node[u, fill=blue!20, draw=blue!80!black] (h_chuc) at (7.0, 0) {HÀNG CHỤC\\$c \times 10$};
  \node[u, fill=orange!20, draw=orange!80!black] (h_dv) at (10.5, 0) {HÀNG ĐƠN VỊ\\$d \times 1$};

  \draw[arr, purple!80!black] (h_nghin) -- (h_tram);
  \draw[arr, teal!80!black] (h_tram) -- (h_chuc);
  \draw[arr, blue!80!black] (h_chuc) -- (h_dv);
\end{tikzpicture}
\end{tcolorbox}

\vspace{-2.5mm}
% --- KHỐI 4: VÍ DỤ & MẸO TRÁNH BẪY ---
\noindent
\begin{minipage}[t]{0.49\textwidth}\vspace{0pt}
\begin{tcolorbox}[
  enhanced,
  colback=green!4!white,
  colframe=green!60!black,
  arc=2mm, boxrule=0.8pt,
  left=2mm, right=2mm, top=1.8mm, bottom=1.8mm,
  title={\footnotesize\bfseries\color{white}\faCheckCircle\ Ví dụ thực hành chuẩn mực}
]
{\footnotesize
\textbf{1. Viết và đọc số:}\\
-- Số gồm $3$ nghìn, $5$ trăm, $0$ chục, $4$ đơn vị viết là: $\mathbf{3504}$. Đọc là: \textit{Ba nghìn năm trăm linh tư}.\\
-- Số gồm $7$ nghìn, $0$ trăm, $8$ chục, $5$ đơn vị viết là: $\mathbf{7085}$. Đọc là: \textit{Bảy nghìn không trăm tám mươi lăm}.\\
\textbf{2. Số liền trước và liền sau:}\\
-- Liền trước số $4560$ là: $4560 - 1 = \mathbf{4559}$.\\
-- Liền sau số $9999$ là: $9999 + 1 = \mathbf{10\,000}$.
}
\end{tcolorbox}
\end{minipage}
\hfill
\begin{minipage}[t]{0.49\textwidth}\vspace{0pt}
\begin{tcolorbox}[
  enhanced,
  colback=red!4!white,
  colframe=red!70!black,
  arc=2mm, boxrule=0.8pt,
  left=2mm, right=2mm, top=1.8mm, bottom=1.8mm,
  title={\footnotesize\bfseries\color{white}\faExclamationTriangle\ Mẹo tránh bẫy trắc nghiệm}
]
{\footnotesize
\begin{itemize}\setlength{\itemsep}{1.5pt}\setlength{\leftskip}{-2.5mm}
  \item \textbf{Bẫy viết thiếu chữ số 0 ở các hàng:}\\
  "Năm nghìn không trăm linh ba", nhiều bạn viết $53$ hoặc $503$ là sai! Số có 4 chữ số phải viết đủ: $\mathbf{5003}$.
  \item \textbf{Quy cách viết số chuẩn quốc gia:}\\
  Dùng khoảng trắng hẹp để phân tách lớp nghìn: viết $10\,000$, TUYỆT ĐỐI không dùng dấu phẩy như tiếng Anh ($10,000$).
\end{itemize}
}
\end{tcolorbox}
\end{minipage}


\end{document}
